\section{Lecture 12: Optimization in Practice}

\subsection{Motivation}

In Lecture 4, we studied gradient-based optimization in an idealized setting.

\medskip

\noindent
However, modern machine learning involves:

\begin{itemize}
    \item large datasets,
    \item high-dimensional parameter spaces,
    \item non-convex objectives.
\end{itemize}

\medskip

\noindent
This raises practical questions:

\medskip

\noindent
\textit{How do we efficiently optimize large-scale models?}

\medskip

\noindent
\textbf{Guiding idea.}  
Efficient optimization relies on stochastic approximations and carefully designed update rules.

\subsection{Stochastic Gradient Descent}

Recall the empirical objective:
\[
L(\theta) = \frac{1}{n} \sum_{i=1}^n \ell_i(\theta).
\]

\medskip

\noindent
\textbf{Gradient descent.}
\[
\theta_{t+1} = \theta_t - \eta \nabla L(\theta_t).
\]

\medskip

\noindent
\textbf{Stochastic gradient descent (SGD).}
\[
\theta_{t+1} = \theta_t - \eta \nabla \ell_i(\theta_t),
\]
where $i$ is sampled randomly.

\medskip

\noindent
\textbf{Properties.}
\begin{itemize}
    \item Low computational cost per update
    \item Noisy gradient estimates
\end{itemize}

\medskip

\noindent
\textbf{Expectation.}
\[
\mathbb{E}[\nabla \ell_i(\theta)] = \nabla L(\theta).
\]

\medskip

\noindent
\textbf{Insight.}  
SGD approximates true gradients using random samples.

\subsection{Mini-Batching and Efficiency}

Instead of a single sample, we often use a mini-batch $B$:

\[
\theta_{t+1} = \theta_t - \eta \frac{1}{|B|} \sum_{i \in B} \nabla \ell_i(\theta_t).
\]

\medskip

\noindent
\textbf{Advantages.}
\begin{itemize}
    \item Reduces variance of gradient estimate
    \item Efficient use of parallel hardware (e.g., GPUs)
\end{itemize}

\medskip

\noindent
\textbf{Tradeoffs.}
\begin{itemize}
    \item Small batches:
    \begin{itemize}
        \item Noisy but fast
    \end{itemize}
    \item Large batches:
    \begin{itemize}
        \item Stable but computationally expensive
    \end{itemize}
\end{itemize}

\medskip

\noindent
\textbf{Insight.}  
Mini-batching balances efficiency and stability.

\subsection{Momentum and Adaptive Methods}

\textbf{Momentum.}

Introduce a velocity variable:
\[
v_{t+1} = \beta v_t + \nabla L(\theta_t),
\]
\[
\theta_{t+1} = \theta_t - \eta v_{t+1}.
\]

\medskip

\noindent
\textbf{Interpretation.}
\begin{itemize}
    \item Accumulates past gradients
    \item Smooths updates
\end{itemize}

\medskip

\noindent
\textbf{Effect.}
\begin{itemize}
    \item Accelerates convergence in consistent directions
    \item Reduces oscillations in ill-conditioned problems
\end{itemize}

\medskip

\noindent
\textbf{Adaptive methods.}

Adjust learning rates per parameter:
\[
\theta_{t+1} = \theta_t - \eta D_t^{-1} \nabla L(\theta_t),
\]
where $D_t$ depends on past gradients.

\medskip

\noindent
\textbf{Examples (conceptually).}
\begin{itemize}
    \item Scale updates based on gradient magnitude
    \item Larger steps for infrequent features
\end{itemize}

\medskip

\noindent
\textbf{Tradeoffs.}
\begin{itemize}
    \item Faster initial progress
    \item Sometimes poorer generalization
\end{itemize}

\medskip

\noindent
\textbf{Insight.}  
Optimization methods influence both convergence and generalization.

\subsection{Practical Convergence Behavior}

\textbf{Non-convexity.}
\begin{itemize}
    \item Many local minima and saddle points
    \item Exact convergence guarantees are limited
\end{itemize}

\medskip

\noindent
\textbf{Empirical observations.}
\begin{itemize}
    \item SGD often finds good solutions
    \item Noise helps escape saddle points
\end{itemize}

\medskip

\noindent
\textbf{Learning rate.}
\begin{itemize}
    \item Critical hyperparameter
    \item Often decreased over time
\end{itemize}

\medskip

\noindent
\textbf{Generalization effects.}
\begin{itemize}
    \item Noisy updates can favor flatter minima
    \item Optimization dynamics act as implicit regularization
\end{itemize}

\medskip

\noindent
\textbf{Stopping criteria.}
\begin{itemize}
    \item Fixed number of iterations
    \item Convergence of loss
    \item Validation performance (early stopping)
\end{itemize}

\medskip

\noindent
\textbf{Insight.}  
Optimization behavior in practice differs significantly from classical theory.

\subsection{A Unifying Perspective}

Optimization in modern machine learning involves:

\begin{itemize}
    \item \textbf{Stochasticity:} approximate gradients
    \item \textbf{Dynamics:} iterative updates in parameter space
    \item \textbf{Adaptivity:} adjusting updates based on history
    \item \textbf{Regularization:} implicit effects of optimization
\end{itemize}

\medskip

\noindent
\textbf{Core idea.}  
Optimization is not just a computational tool, but a fundamental component of learning.

\medskip

\noindent
\textbf{Key insight.}  
The choice of optimization algorithm affects both training efficiency and generalization performance.

\subsection{Outlook}

In this lecture, we studied:

\begin{itemize}
    \item stochastic gradient descent,
    \item mini-batching,
    \item momentum and adaptive methods,
    \item practical convergence behavior.
\end{itemize}

\medskip

\noindent
These ideas are essential for training modern machine learning models.

\medskip

\noindent
In the next lecture, we examine numerical and statistical stability, focusing on how data scaling and noise affect learning.

\medskip

\noindent
\textbf{Guiding principle.}  
Effective learning depends on both the objective and the optimization dynamics used to minimize it.